%==========================================================================================
\section{So, is one parity more important?}\label{sec:which}
\begin{question}
For a harder problem one would like to know whether one parity can be neglected. What do the results say?
\end{question}
The answer depends on which of three questions is being asked.

\paragraph{1. Kinematics: no.} The two parities carry energy with identical weights, Eq.~\eqref{eq:flux}, and they are
exchanged by a $45^\circ$ rotation of the polarization. A plane gravitational wave contains equal amounts of both: its
partial-wave expansion has $\Psi^{\ell,\pm2}_{\rm odd}=\pm i\,\Psi^{\ell,\pm2}_{\rm even}$ for every $\ell$. For \emph{scattering of an
external wave} both parities are equally excited, and dropping one loses half of the signal and all of the
polarization physics of Secs.~\ref{sec:shell}--\ref{sec:time}.

\paragraph{2. Response of the object: comparable in energy, different in phase and in physics.}
\begin{itemize}
\item For a non-absorbing object (the shell) $|S_{\rm even}|=|S_{\rm odd}|=1$. Both parities return exactly the energy they
bring. For a black hole the moduli are equal by isospectrality. With an absorbing (``open'') interior the two
reflectivities differ only at intermediate frequencies, near the barrier top, and coincide at low and high frequency.
\item What differs is the \emph{phase} and therefore the \emph{time structure}: the retardance, the ringing frequencies and
their damping (Secs.~\ref{sec:shell}, \ref{sec:time}).
\item Qualitatively, only the even sector \emph{moves the shell and couples to its matter}: radial displacement,
compression and the equation of state all enter through $\Psi_{\rm even}$. The odd sector sees a universal $\delta$-barrier,
$[\partial_x\Psi]=\Delta\Psi$ with $\Delta=\sqrt{f_R}(1-\sqrt{f_R})/R$, which does not depend on the frequency or on the
equation of state. The matter resonances, their long ringing, and the instability of the static shell exist
\emph{only} in the even sector.
\end{itemize}

\paragraph{3. Sources: even dominates for slow sources.} For a source with internal velocities $v\ll1$, the even
$\ell=2$ mode is the mass quadrupole and the odd $\ell=2$ mode is the current quadrupole \cite{Thorne80}. The latter is
suppressed by one power of $v$ in amplitude. For a circular binary of masses $m_1,m_2$ and total mass $m$, the post-Newtonian
flux gives \cite{Blanchet14}
\begin{equation}
\frac{\dot E_{\rm odd,\ \ell=2}}{\dot E_{\rm even,\ \ell=2}}\simeq\frac{1}{36}\Big(\frac{m_1-m_2}{m}\Big)^2v^2 ,
\end{equation}
of the same order as the even $\ell=3$ (mass-octupole) contribution. A radially infalling or head-on source is
axisymmetric and polar and excites \emph{only} the even sector. Odd parity is dominant only for sources driven by
currents: rotation, toroidal flows, or axial (``$w$-mode'' or $r$-mode-like) oscillations.

\begin{keyidea}
Advice for a harder problem. (i) If the problem is the scattering of an external wave, keep both parities. They carry equal
energy, and the polarization of the outgoing wave is set by their relative phase. (ii) If the waves are generated by slow
matter motion, the even sector dominates and the odd sector is an $O(v^2)$ correction in power. (iii) In either case the
odd sector is the \emph{simpler} and more universal one: for a thin shell it reduces to a $\delta$-barrier. It is a
good first test of the code, while the even sector contains the matter physics and the possible instabilities.
\end{keyidea}
